% Build with XeLaTeX: TinyTeX + DejaVu Sans; the Makefile next door does it.
\documentclass[aspectratio=169,10pt]{beamer}

\usepackage{polyglossia}
\setmainlanguage{english}
\usepackage{fontspec}
\setmainfont{DejaVu Sans}
\setsansfont{DejaVu Sans}
\setmonofont{DejaVu Sans Mono}[Scale=0.85]

\usepackage{tikz}
\usetikzlibrary{positioning, arrows.meta, backgrounds, calc}
\usepackage{booktabs}

% --- the look: calm, no ornament
\usetheme{default}
\usecolortheme{default}
\setbeamertemplate{navigation symbols}{}
\definecolor{ink}{HTML}{1F2933}
\definecolor{accent}{HTML}{1B6B7A}
\definecolor{warm}{HTML}{B4531A}
\definecolor{muted}{HTML}{7B8794}
\definecolor{soft}{HTML}{EEF2F4}
\setbeamercolor{structure}{fg=accent}
\setbeamercolor{normal text}{fg=ink}
\setbeamercolor{frametitle}{fg=ink, bg=}
\setbeamerfont{frametitle}{size=\Large, series=\bfseries}
\setbeamerfont{framesubtitle}{size=\small, series=\mdseries}
\setbeamercolor{framesubtitle}{fg=muted}
\setbeamertemplate{frametitle}{%
  \vspace{0.7em}\usebeamerfont{frametitle}\insertframetitle%
  \ifx\insertframesubtitle\@empty\else
    \\[0.15em]\usebeamerfont{framesubtitle}\color{muted}\insertframesubtitle
  \fi\vspace{0.2em}}
\setbeamertemplate{itemize item}{\color{accent}\raisebox{0.12em}{\rule{0.45em}{0.45em}}}
\setbeamertemplate{itemize subitem}{\color{muted}--}
\setbeamerfont{itemize/enumerate body}{size=\normalsize}
\setbeamerfont{itemize/enumerate subbody}{size=\small}
\setlength{\leftmargini}{1.2em}

% the footer: the page number and a short label
\setbeamertemplate{footline}{%
  \hbox{}\hfill\raisebox{1.2em}{\color{muted}\scriptsize
    ProLecture\quad\insertframenumber/\inserttotalframenumber}\hspace{1.5em}}

\newcommand{\lead}[1]{{\color{warm}\textbf{#1}}}
\newcommand{\num}[1]{{\color{accent}\textbf{#1}}}

\title{ProLecture}
\subtitle{A benchmark of proactive tasks for streaming VLMs}
\date{}

\begin{document}

% ------------------------------------------------------------------ 1
\begin{frame}[plain]
  \vspace{1.2cm}
  {\Huge\bfseries ProLecture}\\[0.5em]
  {\large\color{muted} A benchmark of proactivity that runs into reasoning}\\[2.2em]

  \begin{minipage}{0.9\textwidth}
    \color{ink}
    The model watches a lecture in real time and decides for itself\\
    \lead{when} to answer and \lead{what} to answer.
  \end{minipage}

  \vfill
  {\color{muted}\small The material is 3Blue1Brown lectures. Status: a pilot, 10 samples, the annotation pipeline assembled.}
  \vspace{0.8cm}
\end{frame}

% ------------------------------------------------------------------ 2
\begin{frame}
  \frametitle{The niche}
  \framesubtitle{proactivity where the answer requires a derivation rather than a reading}

  \begin{center}
  \footnotesize
  \begin{tabular}{@{}lll@{}}
    \toprule
    \textbf{Work} & \textbf{What it measures} & \textbf{What is missing} \\
    \midrule
    OVO-Bench & backward / real-time / forward & MCQ, no audio, perception \\
    ProactiveVideoQA & open-ended, the PAUC metric & simple questions ``tell me when X'' \\
    OmniPro & 6 abilities $\times$ 9 subtasks & 84\% audio, yet vision-only keeps up \\
    \bottomrule
  \end{tabular}
  \end{center}

  \vspace{0.8em}
  \begin{itemize}
    \item A lecture has almost \textbf{no perceptual difficulty}: clean rendering, large formulas, clear speech.
    \item What it does have is a logical structure: premise $\rightarrow$ derivation $\rightarrow$ consequence.
    \item So the only way to answer in time is to \lead{understand the content} — not to recognise an object.
  \end{itemize}
\end{frame}

% ------------------------------------------------------------------ 3
\begin{frame}
  \frametitle{Why 3Blue1Brown}
  \framesubtitle{a unique asset: we have the source code of the picture}

  \begin{columns}[T]
    \begin{column}{0.52\textwidth}
      For \num{140} videos both the YouTube link \emph{and} an existing path to the
      manim source of the scene are known.\\[0.6em]
      None of the existing benchmarks has the sources of the video track.
    \end{column}
    \begin{column}{0.44\textwidth}
      \begin{itemize}
        \item an exact ground truth of what is on the screen (the \texttt{Tex} strings with \LaTeX);
        \item semantic boundaries: scene names, the order of the animations;
        \item \lead{a re-render with an injected error} — content guaranteed not to have
              been in pre-training.
      \end{itemize}
    \end{column}
  \end{columns}

  \vspace{1.0em}
  {\color{muted}\small
  The connection ``lesson $\leftrightarrow$ video $\leftrightarrow$ code'' exists only in the frontmatter of the site:
  the \texttt{3b1b/videos} repository itself contains not a single link to YouTube.}
\end{frame}

% ------------------------------------------------------------------ 4
\begin{frame}
  \frametitle{The anatomy of a sample}
  \framesubtitle{the question in one second, the answer in a window, the evidence separately}

  \begin{center}
  \begin{tikzpicture}[x=1.05cm, y=1cm]
    \draw[-{Latex[length=3mm]}, thick, muted] (0,0) -- (12.4,0);
    \node[muted, anchor=west, font=\small] at (12.2,-0.45) {time};

    % the question
    \draw[warm, very thick] (2.2,-0.35) -- (2.2,0.9);
    \node[warm, anchor=south, font=\small\bfseries] at (2.2,0.95) {question\_time};

    % the answer window
    \fill[accent, opacity=0.16] (6.0,-0.32) rectangle (9.2,0.32);
    \draw[accent, thick] (6.0,-0.32) rectangle (9.2,0.32);
    \node[accent, anchor=south, font=\small\bfseries] at (7.6,0.4) {answer window};

    % the distractor
    \fill[warm, opacity=0.12] (3.4,-0.32) rectangle (4.9,0.32);
    \draw[warm, dashed] (3.4,-0.32) rectangle (4.9,0.32);
    \node[warm, anchor=north, font=\scriptsize] at (4.15,-0.42) {distractor};

    % the evidence
    \draw[ink, very thick] (6.3,-1.15) -- (8.6,-1.15);
    \draw[ink] (6.3,-1.25) -- (6.3,-1.05);
    \draw[ink] (8.6,-1.25) -- (8.6,-1.05);
    \node[ink, anchor=north, font=\scriptsize] at (7.45,-1.3) {evidence\_span — where the answer \emph{becomes possible}};

    \node[muted, anchor=north, font=\scriptsize] at (2.2,-0.45) {question issued};
  \end{tikzpicture}
  \end{center}

  \begin{itemize}
    \item The \texttt{temporal\_rel} axis (\texttt{pre} / \texttt{at} / \texttt{post}) is \lead{checked by the validator}:
          it recomputes it from \texttt{question\_time} and \texttt{evidence\_span} and rejects a mismatch.
    \item \texttt{distractor\_windows} — moments where an answer looks apt but is premature.
    \item The confirmation of the reference answer lives separately, in \texttt{provenance}: it may lie later than the window and even outside the video.
  \end{itemize}
\end{frame}

% ------------------------------------------------------------------ 5
\begin{frame}
  \frametitle{What else is in a sample}
  \framesubtitle{what the answer is judged by, and what in the annotation is measured rather than declared}

  \begin{columns}[T]
    \begin{column}{0.47\textwidth}
      {\bfseries In every window}\\[0.4em]
      {\small
      \texttt{response} — the reference answer\\[0.2em]
      \texttt{rubric} — what the judge scores against\\[0.2em]
      \texttt{trigger\_modality} — visual / speech / sound\\[0.2em]
      \texttt{depends\_on} — a reference to a previous window\\[0.5em]
      The rubric is what replaces the MCQ. From \texttt{q01}:
      \begin{itemize}\setlength{\itemsep}{0pt}
        \item[] ``the number 2.5 is named''
        \item[] ``if a computation is shown, it is weighted, with 1/2, 1/4, 1/8, 1/8''
        \item[] \lead{``do NOT credit 1.75 — that is the original code, not the swapped one''}
      \end{itemize}}
    \end{column}
    \begin{column}{0.47\textwidth}
      {\bfseries \texttt{validation} — the results of the measurements}\\[0.4em]
      {\small
      One flag per probe, each filled in by a \emph{model's answer} rather than by an opinion:\\[0.3em]
      \texttt{solvable} — is it solvable at all\\[0.2em]
      \texttt{blind\_solved} — solvable without the video\\[0.2em]
      \texttt{textonly\_solved} — the transcript is enough\\[0.2em]
      \texttt{visiononly\_solved} — the frames are enough\\[0.2em]
      \texttt{frozenframe\_solved} — one frame is enough\\[0.6em]
      \texttt{audio\_dependency} is derived from these; the annotator's opinion lives separately,
      in \texttt{audio\_dependency\_claim}, and influences no conclusion.\\[0.4em]
      \texttt{status}: \texttt{draft} $\rightarrow$ \texttt{validated} / \texttt{rejected}.}
    \end{column}
  \end{columns}
\end{frame}

% ------------------------------------------------------------------ 6
\begin{frame}
  \frametitle{The three taxonomy axes}
  \framesubtitle{every sample is a point in the cube (act, relation, level of reasoning)}

  \begin{columns}[T]
    \begin{column}{0.33\textwidth}
      {\bfseries A. The type of act}\\[0.4em]
      {\small
      \texttt{alert} — an alert\\
      \texttt{monitor} — a state\\
      \texttt{ground} — grounding\\
      \texttt{narrate} — narration\\
      \texttt{predict} — prediction\\
      \texttt{verify} — verification\\
      \texttt{counterfact} — a counterfactual}
    \end{column}
    \begin{column}{0.33\textwidth}
      {\bfseries B. The relation to the trigger}\\[0.4em]
      {\small
      \texttt{pre} — the evidence comes \emph{after} the question: \lead{the ability to wait}\\[0.3em]
      \texttt{at} — here and now\\[0.3em]
      \texttt{post} — an answer from memory\\[0.5em]
      \color{muted} One question in two regimes gives a controlled pair.}
    \end{column}
    \begin{column}{0.33\textwidth}
      {\bfseries C. The level of reasoning}\\[0.4em]
      {\small
      \texttt{L1} perception — read off\\[0.3em]
      \texttt{L2} integration — connect two moments\\[0.3em]
      \texttt{L3} derivation — carry out reasoning that is not in the video\\[0.5em]
      \color{muted} The criterion for L3: the blind baseline does not solve it.}
    \end{column}
  \end{columns}

  \vspace{0.9em}
  {\color{muted}\small Axis C is what OmniPro does not have: it is precisely along it that its models fail to differ.}
\end{frame}

% ------------------------------------------------------------------ 7
\begin{frame}
  \frametitle{The principles the design grew out of}
  \framesubtitle{each point closes a failure that sank one of the works reviewed}

  \begin{itemize}
    \item \lead{Labels are set by measurement, not by declaration.} \texttt{audio\_dependency}
          is derived from the modality probes; the annotator's opinion lives in a separate
          field and influences no conclusion.
    \item \lead{Generator $\neq$ judge $\neq$ participant} — a rule \emph{for the measurement phase}.
          At present the proposer, the reviewers and the coordinator are one model, and no judge
          or participant has been run yet. The judge may be a pure text LLM: it reads the rubric
          and the answer, not the video.
    \item \lead{Open answers with a rubric, no MCQs.} A skewed answer key breaks a benchmark quietly.
    \item \lead{Mandatory dummy baselines:} blind, text-only, vision-only, frozen-frame,
          ``answer immediately'', ``answer at the end'', ``stay silent''.
    \item \lead{No quotas} of the form ``exactly N samples per task'': however many passed validation is however many there are.
    \item \lead{Contamination is measured, not denied:} a 2026 pilot video and a 2017 control
          one, so the blind-baseline gap becomes a number.
  \end{itemize}
\end{frame}

% ------------------------------------------------------------------ 8
\begin{frame}
  \frametitle{How this works on live samples}
  \framesubtitle{one and the same blind check, opposite outcomes}

  \begin{columns}[T]
    \begin{column}{0.47\textwidth}
      {\color{accent}\bfseries q01 — passed the check}\\[0.4em]
      {\small
      \emph{``As soon as a table of code words of different lengths with their probabilities
      first appears, compute what the average number of bits becomes if `up' and `left'
      are swapped.''}\\[0.5em]
      Question at 250 s, table at 314 s; the answer of 2.5 bits is not in the video.\\[0.5em]
      The semi-blind reviewer \emph{named} 2.5 — and was not credited: the answer opened with
      ``I cannot name an exact number'', then a conditional guess and a branch with 2.0.
      \lead{The problem was guessed, the answer was not found.}}
    \end{column}
    \begin{column}{0.47\textwidth}
      {\color{warm}\bfseries q08 — rejected}\\[0.4em]
      {\small
      \emph{``Say before the author tells it himself: what source of letter probabilities
      did Shannon turn to, and why does he not have the frequency-table problem?''}\\[0.5em]
      Before any viewing the reviewer produced \textbf{both} halves of the answer: a live
      native speaker, and the growth of the number of contexts as $27^n$. Three rubric items
      out of three, without a single assumption about a frame.\\[0.5em]
      \lead{\texttt{blind\_solved = true}}, \texttt{status: rejected}; the level was lowered
      \texttt{L3} $\rightarrow$ \texttt{L1} along the way.}
    \end{column}
  \end{columns}

  \vspace{0.2em}
  {\color{muted}\footnotesize The rule: leaning on a commonly known fact is contamination; guessing the problem is not.}
\end{frame}

% ------------------------------------------------------------------ 9
\begin{frame}
  \frametitle{The pipeline}
  \framesubtitle{from a video to a sample that has been through review}

  \begin{center}
  \begin{tikzpicture}[
      node distance=0.28cm,
      box/.style={draw=accent, fill=soft, rounded corners=2pt, inner sep=4pt,
                  align=center, font=\scriptsize, minimum height=1.05cm, text width=1.6cm},
      lab/.style={font=\scriptsize, color=muted, align=center},
      ar/.style={-{Latex[length=2mm]}, muted, thick}]
    \node[box] (a) {video\\+ source};
    \node[box, right=of a] (b) {bundle:\\transcript,\\frames, formulas};
    \node[box, right=of b] (c) {proposer\\(agent)};
    \node[box, right=of c] (d) {validator\\(schema)};
    \node[box, right=of d] (e) {reviewers\\and probes};
    \node[box, right=of e] (f) {human\\acceptance};
    \draw[ar] (a) -- (b); \draw[ar] (b) -- (c); \draw[ar] (c) -- (d);
    \draw[ar] (d) -- (e); \draw[ar] (e) -- (f);
    \draw[ar] (e.south) .. controls +(0,-1.0) and +(0,-1.0) .. (d.south);
    \node[lab] at ($(d.south)!0.5!(e.south)$) [below=0.95cm] {not PASS — the same question goes to fresh subagents};
  \end{tikzpicture}
  \end{center}

  \vspace{0.4em}
  \begin{itemize}
    \item \textbf{The validator} plays the role of tests: non-overlapping windows, references only backwards,
          a distractor not coinciding with a trigger, \texttt{temporal\_rel} agreeing with the evidence.
          It returns code 1 on any error.
    \item \textbf{The reviewers and the probes} are the next two slides: three roles check the annotation
          from different sides, four probes fill in the \texttt{validation} flags by measurement.
    \item \textbf{The round repeats} until every verdict is \texttt{PASS};
          a repeat round goes to fresh subagents that have not seen the previous discussions.
  \end{itemize}
\end{frame}

% ------------------------------------------------------------------ 10
\begin{frame}
  \frametitle{Three reviewers, blind to each other}
  \framesubtitle{they run in parallel; none of them edits the files}

  \begin{itemize}
    \item \lead{The semi-blind one.} Sees the material \emph{up to} \texttt{t\_end} but does not know
          \texttt{t\_start} — and names for themselves the moment the answer became possible.
          A checker who knows the answer ``recognises'' it on a half-drawn frame and keeps moving
          the left boundary of the window leftwards; this one does not have that knowledge.
          \begin{itemize}
            \item[] {\small They separately record \texttt{prior\_answer} — their own answer \textbf{before} viewing,
                    verbatim. It is what the coordinator derives \texttt{blind\_solved} from.}
          \end{itemize}
    \item \lead{The sample reviewer.} Everything relating to the question as a whole: the wording,
          the rubric, the decoys, the reference answer, the links between windows. The window boundaries are not their work.
    \item \lead{The per-window reviewer, one per window.} The role was separated not for symmetry:
          with a single reviewer on a question with nine windows, the frame budget was enough
          for the first, and after that they made things up.
  \end{itemize}

  \vspace{0.5em}
  {\color{muted}\small Each returns a structured review against the contract
  (\texttt{schema/review.py}); the coordinator applies the edits, so it is visible who found what.}
\end{frame}

% ------------------------------------------------------------------ 11
\begin{frame}
  \frametitle{Four probes: the label is measured}
  \framesubtitle{each is given exactly one channel, and we see whether it answers}

  \begin{center}
  \small
  \begin{tabular}{@{}llll@{}}
    \toprule
    \textbf{Probe} & \textbf{Material} & \textbf{Frames} & \textbf{Fills in} \\
    \midrule
    \texttt{text}   & transcript up to \texttt{t\_end} & 0        & \texttt{textonly\_solved} \\
    \texttt{frame}  & one frame mid-window             & 1        & \texttt{frozenframe\_solved} \\
    \texttt{vision} & frames up to \texttt{t\_end}     & up to 40 & \texttt{visiononly\_solved} \\
    \texttt{full}   & transcript and frames            & up to 40 & \texttt{solvable} \\
    \bottomrule
  \end{tabular}
  \end{center}

  \vspace{0.6em}
  {\small
  \texttt{solvable = textonly OR visiononly OR} the \texttt{full} answer; \texttt{solvable = false}
  is a \texttt{CRITICAL\_FAIL} — the question is not solvable even from the full material.\\[0.3em]
  \texttt{visiononly\_solved} together with \texttt{solvable} give \lead{\texttt{audio\_dependency}}
  — this label is never set by hand.\\[0.3em]
  \texttt{frozenframe\_solved = true} at level \texttt{L2}/\texttt{L3} means the level is inflated.}

  \vspace{0.5em}
  {\color{muted}\small A probe's isolation rests on three things: the prompt \emph{asks} it not to reach into
  another channel, the script \emph{does not let} two modes share a folder, and the probe must
  confess in \texttt{violations}.}
\end{frame}

% ------------------------------------------------------------------ 12
\begin{frame}
  \frametitle{Tools that did not exist}
  \framesubtitle{two places where what was available did not fit}

  \begin{columns}[T]
    \begin{column}{0.48\textwidth}
      \textbf{Our own scene change detector}\\[0.5em]
      {\small
      The off-the-shelf detectors compare a frame with the previous one — while manim's typical
      transition is a smooth travel over 2--4 seconds, where neighbouring frames barely differ.\\[0.5em]
      We compute the \lead{accumulated drift} from the last captured frame and capture
      once the animation has settled.\\[0.5em]
      One threshold per domain: \num{13.5} / \num{12.9} / \num{10.1} frames per minute
      on the videos from 2026 and 2017.\\[0.3em]
      On the control stretch: \num{18} states out of 18 against \num{10} for the uniform grid.}
    \end{column}
    \begin{column}{0.48\textwidth}
      \textbf{A dense description of the video track}\\[0.5em]
      {\small
      A document like subtitles, only with what is happening on the screen instead of utterances.
      The annotator reads text instead of looking through dozens of frames.\\[0.5em]
      Two frames per second, windows of 5 seconds with an overlap, strict JSON against a schema
      (guided decoding on the server side).\\[0.5em]
      The pilot: \num{3878} frames, \num{485} fragments, \num{20} minutes on two A100s,
      \num{1939} seconds out of 1939 covered by events.\\[0.3em]
      \color{muted} The describing model is not shown the subtitles: it describes only the frames.
      Otherwise the description could not be given to the vision-only probe instead of frames.}
    \end{column}
  \end{columns}
\end{frame}

% ------------------------------------------------------------------ 13
\begin{frame}
  \frametitle{Where we are now}

  \begin{columns}[T]
    \begin{column}{0.46\textwidth}
      \textbf{Done}
      \begin{itemize}
        \item[] {\small
        the data pipeline: mapping $\rightarrow$ download $\rightarrow$ bundle;\\[0.3em]
        the sample schema and the semantic validator;\\[0.3em]
        the pilot: \num{10} samples, \num{24} triggers and \num{12} distractors,\\ zero validation errors;\\[0.3em]
        the prompts of the proposer, three reviewers and four probes;\\[0.3em]
        the scene detector and the description of the video track.}
      \end{itemize}
    \end{column}
    \begin{column}{0.46\textwidth}
      \textbf{Next}
      \begin{itemize}
        \item[] {\small
        run the dummy baselines and see that they fail;\\[0.3em]
        measure the blind gap on the 2017 control video;\\[0.3em]
        the streaming harness and an LLM judge against the rubric;\\[0.3em]
        expanding the annotation: the share of \texttt{L3} must grow;\\[0.3em]
        optionally — re-rendering scenes with an injected error.}
      \end{itemize}
    \end{column}
  \end{columns}

  \vspace{1.2em}
  \begin{center}
    {\color{muted}\small The open assumptions are written down separately and in Limitations rather than swept away:
    the leak of the distractor's explanation, the high difficulty of the pilot, the small volume.}
  \end{center}
\end{frame}

% ------------------------------------------------------------------ 14
\begin{frame}
  \frametitle{The metric}
  \framesubtitle{designed but not yet implemented: groundwork for the measurement phase}

  \begin{columns}[T]
    \begin{column}{0.48\textwidth}
      \textbf{PAUC.} On each window a polyline ``time $\rightarrow$ quality of the accumulated
      answers'' is built, and the quality $s \in \{0,1,2\}$ is assigned by an LLM judge against the rubric.

      \vspace{0.6em}
      The starting point $(t_{\text{start}}, 0.5)$: silence is better than a confident wrong answer.

      \vspace{0.6em}
      \lead{A consequence we report explicitly:} the strategy ``always stay silent'' gets
      \num{0.25} of the maximum. That is the floor of the scale, not a zero.
    \end{column}
    \begin{column}{0.48\textwidth}
      \textbf{Mandatory alongside:}
      \begin{itemize}
        \item[] {\small
        \textbf{False-Alarm Rate} — answers outside the windows\\[0.25em]
        \textbf{Trap rate} — hits in the distractors, separately from the FAR\\[0.25em]
        \textbf{Latency} — from $t_{\text{start}}$, as a distribution\\[0.25em]
        \textbf{Silence rate}, \textbf{Trigger Recall}\\[0.25em]
        \textbf{Responses-per-sample} — otherwise the spam strategy is invisible}
      \end{itemize}
      \vspace{0.4em}
      {\color{muted}\small Aggregation over samples, not over windows: otherwise the result is a hostage
      of the single simplest sample, the one with nine triggers.}
    \end{column}
  \end{columns}
\end{frame}

\end{document}
